\documentclass[11pt,a4paper]{article}
\usepackage[utf8]{inputenc}
\usepackage[T1]{fontenc}
\usepackage[spanish]{babel}
\usepackage[margin=2.5cm]{geometry}
\usepackage{booktabs}
\usepackage{array}
\usepackage{tabularx}
\usepackage{listings}
\usepackage{xcolor}
\usepackage{enumitem}
\usepackage{parskip}
\usepackage{fancyhdr}
\usepackage{amsmath}
\usepackage{amssymb}

\definecolor{codegray}{RGB}{245,245,245}
\definecolor{codeblue}{RGB}{0,60,180}
\definecolor{codegreen}{RGB}{0,120,0}
\definecolor{codepurple}{RGB}{120,0,120}

\lstdefinestyle{sql}{
  language=SQL,
  backgroundcolor=\color{codegray},
  basicstyle=\ttfamily\small,
  keywordstyle=\color{codeblue}\bfseries,
  commentstyle=\color{codegreen},
  stringstyle=\color{codepurple},
  breaklines=true,
  frame=single,
  framerule=0.4pt,
  xleftmargin=8pt,
  xrightmargin=8pt,
  numbers=none,
  tabsize=2,
  showstringspaces=false,
  morekeywords={FETCH,FIRST,ROWS,ONLY,OFFSET,OVER,PARTITION,DENSE_RANK,ROW_NUMBER,NTILE,LAG,LEAD,FIRST_VALUE,LAST_VALUE,INTERSECT,EXCEPT}
}
\lstset{style=sql}

\pagestyle{fancy}
\fancyhf{}
\rhead{\textcolor{gray}{\small TD7: Ingeniería de Datos — UTDT}}
\lhead{\textcolor{gray}{\small Primer Parcial Modelo 3}}
\rfoot{\thepage}
\renewcommand{\headrulewidth}{0.4pt}

\setlist[enumerate,1]{label=\textbf{\alph*)}}

\begin{document}

\begin{center}
  {\LARGE\bfseries TD7 — Ingeniería de Datos}\\[6pt]
  {\large Universidad Torcuato Di Tella}\\[4pt]
  {\Large\bfseries Primer Parcial — Modelo 3}\\[4pt]
  {\normalsize Duración: 2 horas \quad|\quad Puntaje total: 100 puntos}
\end{center}
\noindent\rule{\linewidth}{0.8pt}
\vspace{2pt}

\noindent\textbf{Instrucciones:} Responda todos los ejercicios. Justifique cada respuesta. No está permitido el uso de material de consulta ni dispositivos electrónicos. Indique claramente el número de ejercicio y subejercicio.

\vspace{6pt}
\noindent\rule{\linewidth}{0.4pt}

% ─────────────────────────────────────────────────────────────
\section*{Ejercicio 1 — Modelo Entidad-Interrelación y Modelo Relacional \hfill\normalfont(30 puntos)}

Una agencia de viajes quiere modelar su sistema de paquetes y reservas. Las reglas son:

\begin{itemize}
  \item Los \textbf{destinos} tienen un código único, ciudad, país y descripción.
  \item Los \textbf{paquetes turísticos} tienen un código, nombre, precio base y duración en días. Un paquete incluye \emph{al menos un} destino; un destino puede estar en varios paquetes. El orden de visita de los destinos dentro de un paquete debe quedar registrado.
  \item Los \textbf{clientes} se identifican por DNI y tienen nombre, apellido, email y teléfono.
  \item Una \textbf{reserva} tiene un número único, fecha de reserva y fecha de viaje. La realiza exactamente un cliente para exactamente un paquete. Un cliente puede tener varias reservas, pero un paquete puede no tener ninguna reserva aún.
  \item Los \textbf{guías de turismo} tienen un legajo, nombre y los idiomas que hablan (atributo multivaluado). Un guía puede estar asignado a varias reservas; cada reserva tiene exactamente un guía asignado.
\end{itemize}

\begin{enumerate}
  \item (20 pts) Dibuje el DER completo. Indique entidades, atributos (marcando el identificador y los atributos multivaluados), interrelaciones, cardinalidades y participaciones.
  \vspace{10cm}
  \item (10 pts) Traduzca al modelo relacional indicando nombre de tabla, atributos, \underline{PK} y \textit{FK} de cada tabla. Preste especial atención al atributo multivaluado.
\end{enumerate}

\newpage

% ─────────────────────────────────────────────────────────────
\section*{Ejercicio 2 — Álgebra Relacional \hfill\normalfont(25 puntos)}

Use el esquema del Mundial 2022:

\begin{center}
\begin{tabular}{ll}
\toprule
\textbf{Relación} & \textbf{Atributos} \\
\midrule
\texttt{Players} & \texttt{id, name, birth\_year, playing\_position, national\_team} \\
\texttt{Matches}  & \texttt{id, home, away, datetime, stage} \\
\texttt{Stages}   & \texttt{id, name} \\
\bottomrule
\end{tabular}
\end{center}

\begin{enumerate}
  \item (4 pts) Obtener \texttt{id} y \texttt{name} de todos los jugadores que se desempeñan como arqueros (\texttt{playing\_position = 'GK'}).

  \vspace{2cm}

  \item (4 pts) Obtener los \texttt{name} de los jugadores de Argentina \textbf{o} Brasil. Resuelva usando unión de conjuntos.

  \vspace{2.5cm}

  \item (5 pts) Obtener los \texttt{name} de los jugadores de Argentina que nacieron antes de 1990.

  \vspace{2.5cm}

  \item (6 pts) Para los partidos donde Alemania (\texttt{home = 'GER'}) fue local, obtener el \texttt{datetime} del partido y el \texttt{name} de la etapa correspondiente.

  \vspace{3cm}

  \item (6 pts) Obtener los \texttt{name} de los jugadores que pertenecen a selecciones que jugaron \textbf{al menos un partido como local}. \emph{(Sugerencia: use renombramiento $\rho$ y junta.)}

  \vspace{3.5cm}
\end{enumerate}

\newpage

% ─────────────────────────────────────────────────────────────
\section*{Ejercicio 3 — SQL \hfill\normalfont(45 puntos)}

Use la base de datos \textbf{Chinook} (mismo esquema de los modelos anteriores).

\begin{enumerate}

  \item (7 pts) Listar el \texttt{Title} de todos los álbumes junto con el \texttt{Name} del artista. Incluir únicamente álbumes que tengan artista registrado. Ordenar por nombre del artista y luego por título del álbum.

  \vspace{4cm}

  \item (8 pts) Para cada tipo de medio (\texttt{MediaType}), mostrar su nombre y la cantidad de canciones que tienen una duración mayor a 5 minutos (300\,000 milisegundos). Incluir los tipos de medio aunque no tengan ninguna canción larga (mostrando 0).

  \vspace{4.5cm}

  \item (8 pts) Listar los artistas que tienen al menos un álbum pero para los cuales \textbf{ninguna canción} de ese álbum tiene como \texttt{Composer} el nombre del artista. Usar subconsultas con \texttt{EXISTS} y \texttt{NOT EXISTS}.

  \vspace{5cm}

  \item (8 pts) Encontrar los clientes que hayan comprado canciones de \textbf{más de 3 géneros distintos} (a través de \texttt{Invoice} $\to$ \texttt{InvoiceLine} $\to$ \texttt{Track}). Mostrar \texttt{FirstName}, \texttt{LastName} y la cantidad de géneros distintos.

  \vspace{5cm}

  \item (14 pts) Para cada factura (\texttt{Invoice}), mostrar el \texttt{InvoiceId}, el \texttt{FirstName} y \texttt{LastName} del cliente, el \texttt{Total} de esa factura y el \textbf{monto acumulado} de todas las facturas del mismo cliente hasta esa fecha (inclusive), ordenadas por \texttt{InvoiceDate}. Usar una función de ventana con \texttt{SUM} y \texttt{ROWS BETWEEN UNBOUNDED PRECEDING AND CURRENT ROW}.

  \vspace{5.5cm}

\end{enumerate}

\end{document}
